\documentclass[11pt]{article}
\usepackage{amsmath,amssymb,amsthm,hyperref}
\newtheorem{theorem}{Theorem}
\newtheorem{lemma}[theorem]{Lemma}
\newtheorem{proposition}[theorem]{Proposition}
\newcommand{\R}{\mathbb R}
\newcommand{\N}{\mathbb N}
\newcommand{\1}{\mathbf 1}
\title{A Superlinear One-Hot Decoder from Interval Orders}
\author{Research checkpoint, 23 November 2026}
\date{}
\begin{document}
\maketitle

This note proves a local sign-only decoder with a logarithmic saving. It does
not prove an upper bound for arbitrary truth tables on the full Boolean cube.
The distinction is essential: all denominators below are admissible on the
whole cube, but the arbitrary labels are prescribed only on its one-hot
product subset. No model definition is changed.

\begin{theorem}[Local decoder]
There is an absolute constant $C$ such that for all sufficiently large $m$
and every positive integer $s$, every sign table
$f:\{1,\ldots,m\}\times\{1,\ldots,s\}\to\{-1,1\}$ is represented, with
strict signs, on
\[
 \{(e_i,e_j):1\leq i\leq m,\ 1\leq j\leq s\}\subseteq\{0,1\}^{m+s}
\]
by at most $Cm/\log m$ admissible affine-ratio heads and a constant.
The resulting score can also be chosen nonzero at every vertex of the full
Boolean cube. Every denominator has positive constant and strictly positive slopes on all
$m+s$ original input variables. The same fixed row coefficients and numerator coefficients work for
all columns and all sign tables; only the column denominator coefficients
depend on $f$.
\end{theorem}

The combinatorial source is the interval-order shattering mechanism in
Gillibert, Lachmann, and M\"ullner,
\emph{The VC-Dimension of Axis-Parallel Boxes on the Torus},
Journal of Complexity 68 (2022), 101600,
\url{https://arxiv.org/abs/2004.13861}.
Their sharper result is asymptotic to $d\log_2 d$ points in $d$ coordinates.
For completeness, a weaker constant-factor bound sufficient here follows
from the elementary argument below. This note uses nonwrapping intervals,
so no periodic activation or periodic arithmetic is being introduced.

\begin{lemma}[Extraction matrix]
For every sufficiently large integer $d$, put $c=\lfloor d/2\rfloor$ and
$k=\lfloor\sqrt d\rfloor$. There is a matrix
$E\in\{1,\ldots,k\}^{c\times d}$ with this property: for every word
$(a_1,\ldots,a_c)\in\{1,\ldots,k\}^c$ there are distinct columns
$j_1,\ldots,j_c$ satisfying $E_{i,j_i}=a_i$.
\end{lemma}
\begin{proof}
Choose all matrix entries independently and uniformly from the alphabet.
For a fixed requested word let
$X_i=\{j:E_{i,j}=a_i\}$. Hall's marriage theorem says that distinct
representatives exist if every $U\subseteq\{1,\ldots,c\}$ satisfies
$|\bigcup_{i\in U}X_i|\geq |U|$.
Consequently, failure for any word implies that, for some nonempty set
$U$ of size $u\leq c$ and some set $V$ of $u-1$ columns, every row in $U$
has an alphabet symbol absent from all columns outside $V$.
For fixed $U,V$, independence between rows and a union bound over the symbol
in each row give probability at most
\[
 \left[k(1-1/k)^{d-u+1}\right]^u
 \leq \left[k\exp(-d/(2k))\right]^u.
\]
There are at most $d^u$ choices for $U$ and at most $d^u$ choices for $V$.
Thus the probability of any failure is at most
\[
 \sum_{u=1}^c\left[d^2k\exp(-d/(2k))\right]^u.
\]
Since $k\leq\sqrt d$, the bracket is at most
$d^{5/2}\exp(-\sqrt d/2)$, which tends to zero.
For all sufficiently large $d$ it is less than $1/3$, and the displayed
finite geometric sum is less than $1/2$. Some matrix therefore has the
asserted property. This argument ranges over all words simultaneously;
there is no target-dependent choice of matrix.
\end{proof}

\begin{lemma}[Fixed interval orders]
For all sufficiently large $d$, there are $M\geq (d\log_2 d)/16$ points and
$d$ fixed linear orders on those points such that every subset is the union
of one interval in each order, permitting empty intervals.
\end{lemma}
\begin{proof}
Use the extraction matrix above and put $q=\lfloor\log_2 k\rfloor$.
Index the points by pairs $(i,b)$ with $1\leq i\leq c$ and $1\leq b\leq q$,
so $M=cq$. Associate the first $2^q$ alphabet symbols with all subsets
$A\subseteq\{1,\ldots,q\}$; assign arbitrary subsets to remaining symbols.
For coordinate $h\in\{1,\ldots,d\}$, order the $c$ groups by $i$.
Within group $i$, put first those points $(i,b)$ for which $b$ lies in the
subset associated to $E_{i,h}$, and then all remaining points. Break ties by
$b$. This defines a fixed permutation of the $M$ points.

Given a desired subset $S$ of the points, let $a_i$ be the symbol associated
to $\{b:(i,b)\in S\}$. Choose distinct columns $h_i$ by extraction. In
order $h_i$, the desired subset of group $i$ is a consecutive initial block
of that group. Select exactly this block as the interval for coordinate
$h_i$, or select the empty interval if the block is empty. Select the empty
interval for every unused coordinate. These intervals are contained in
separate groups, so their union is exactly $S$.
Finally, for sufficiently large $d$,
$c\geq d/3$ and $q\geq (\log_2 d)/5$, giving the weaker stated constant.
\end{proof}

\begin{lemma}[Positive rational prefix decoder]
Let $R>1$, let $r\geq1$ and $\ell\geq0$ be integers, and define
\[
 K_R(r,\ell)=\frac{R^{2\ell+1}}{1+R^{2\ell+1}+R^{2r}},\qquad
 J_R(r,\ell)=\frac{R^{2r}}{1+R^{2\ell+1}+R^{2r}}.
\]
Then
\[
 0<K_R(r,\ell)<1,\qquad
 \left|K_R(r,\ell)-\1_{r\leq\ell}\right|\leq\frac2R,\qquad
 \left|J_R(r,\ell)-\1_{r>\ell}\right|\leq\frac2R.
\]
\end{lemma}
\begin{proof}
The positivity claims follow from positivity of all powers.
If $r\leq\ell$, then $R^{2r}\leq R^{2\ell}$ and
$1\leq R^{2\ell}$. Writing $T=R^{2\ell+1}$ and $B=R^{2r}$ gives
\[
 1-K_R(r,\ell)=\frac{1+B}{1+T+B}
 \leq\frac{1+B}{T}\leq\frac{2R^{2\ell}}{R^{2\ell+1}}=\frac2R.
\]
If $r>\ell$, then $R^{2r}\geq R^{2\ell+2}$, hence
\[
 K_R(r,\ell)\leq\frac{R^{2\ell+1}}{R^{2r}}\leq\frac1R\leq\frac2R.
\]
For the suffix expression $J_R$, if $r\leq\ell$ then $J_R\leq B/T\leq1/R$.
If $r>\ell$, then $1-J_R\leq(1+T)/B\leq2/R$, because
$B\geq R^{2\ell+2}$ and $1\leq R^{2\ell+1}=T$.
\end{proof}

\begin{proof}[Proof of the local decoder theorem]
First suppose that there are $M$ rows and fixed interval orders as above.
Let $\pi_h(i)\in\{1,\ldots,M\}$ be the rank of row $i$ in order $h$.
For every column $j$, express $\{i:f(i,j)=1\}$ as a union of intervals
in the fixed orders. Write the $h$-th interval as
\[
 \{i:L_{hj}<\pi_h(i)\leq U_{hj}\},\qquad
 0\leq L_{hj}\leq U_{hj}\leq M.
\]
Empty intervals have $L_{hj}=U_{hj}$. Put $R=32d+1$, and define
\[
 b_{hi}=R^{2\pi_h(i)},\qquad
 t^+_{hj}=R^{2U_{hj}+1},\qquad t^-_{hj}=R^{2L_{hj}+1}.
\]
For arbitrary Boolean vectors $u\in\{0,1\}^M$ and
$v\in\{0,1\}^s$, set
\[
 B_h(u)=\sum_i b_{hi}u_i,\qquad
 T_h^\pm(v)=\sum_j t^\pm_{hj}v_j,
\]
\[
 F_h^\pm(u,v)=\frac{B_h(u)}{1+B_h(u)+T_h^\pm(v)},
 \qquad
 S(u,v)=-\frac12+\sum_{h=1}^d(F_h^-(u,v)-F_h^+(u,v)).
\]
Each $F_h^\pm$ is an affine ratio with positive constant denominator and
strictly positive denominator coefficients on every row and column bit.
Its negative is also such an atom, by changing the numerator sign only.
The repository's exact positive-affine-ratio realization therefore gives
$2d$ actual heads.

At $(u,v)=(e_i,e_j)$, each $F_h^\pm$ is the suffix decoder above.
Thus $F_h^--F_h^+$ differs from its interval indicator by at most $4/R$,
and the full sum differs from the integer number of intervals containing
$i$ by at most $4d/R<1/8$. If $f(i,j)=-1$ that integer is zero, and
$S(e_i,e_j)<-3/8$. If $f(i,j)=1$ the integer is at least one, and
$S(e_i,e_j)>3/8$. All prescribed signs are strict. To make every output on the full cube
strict as well, choose $\delta\in(-1/16,1/16)$ avoiding the finitely many
values $-S(u,v)$ and replace $S$ by $S+\delta$. Such a choice exists because
a finite set cannot contain a nonempty real interval. This preserves the
prescribed signs and all denominator properties. It merely chooses the
labels of the extension outside the prescribed one-hot product.

For an arbitrary sufficiently large row count $m$, choose
$d=\lceil A m/\log_2 m\rceil$ with a sufficiently large absolute $A$.
The bound $M\geq d\log_2(d)/16$ is eventually at least $m$.
Restrict the orderings to any $m$ rows. The interval-union property is
preserved by restriction; assign ranks again if desired. The resulting
head count $2d$ is $O(m/\log m)$, uniformly in $s$ and $f$.
\end{proof}

\begin{proposition}[Exact behavior off the one-hot product]
For every Boolean $u,v$ and every $h$, put $D_h=1+B_h(u)$.
Then
\[
 F_h^--F_h^+=
 \frac{B_h(u)\bigl(T_h^+(v)-T_h^-(v)\bigr)}
 {(D_h+T_h^+(v))(D_h+T_h^-(v))}\geq0.
\]
If $u=0$ or $v=0$, this difference is exactly zero, so the complete score
before the final bias perturbation is $-1/2$ on both pure faces. For general nonempty $u,v$, the same formula
applies without approximation: the module compares sums of the fixed row
powers with sums of the column-dependent endpoint powers. It does not
select one row and one column.
\end{proposition}
\begin{proof}
Subtract the two fractions over their common positive denominator.
Since $U_{hj}\geq L_{hj}$, one has $t^+_{hj}\geq t^-_{hj}$, hence
$T_h^+(v)\geq T_h^-(v)$. If $u=0$ then $B_h(u)=0$; if $v=0$ then
$T_h^+(v)=T_h^-(v)=0$. Both vanishing claims are exact.
\end{proof}

This nonnegativity permits monotone union assembly only when the intended
positive regions are already compatible. It does not establish that arbitrary
labels on multiple translated cells can be combined. In particular,
complementing selected input coordinates changes the signs of their
denominator slopes; positivity on a cell alone would not repair the model's
whole-cube admissibility requirement. The global transfer remains open.


\begin{proposition}[Exact dual-rail denominator lift]
Suppose $D(z)=b+\sum_{i=1}^r a_i z_i$ is positive at every vertex of
$\{0,1\}^r$, allowing arbitrary signs for the logical slopes $a_i$.
There is a positive-constant, strictly positive-slope affine denominator
$\widetilde D$ on $\{0,1\}^{2r}$ such that
$\widetilde D(z,1-z)=D(z)$ for every logical input $z$.
Consequently every affine numerator divided by $D$ is exactly one actual
admissible attention atom on the complementary-pair subset.
\end{proposition}
\begin{proof}
Put $\mu=b-\sum_i\max(-a_i,0)$. This is the minimum of $D$ on the
logical cube, attained by setting $z_i=1$ precisely when $a_i<0$; hence
$\mu>0$. Choose $\varepsilon=\mu/(2(r+1))$ and define
\[
 \widetilde b=b-\sum_i\max(-a_i,0)-r\varepsilon>0,\qquad
 w_i=\max(a_i,0)+\varepsilon>0,\qquad
 w_i'=\max(-a_i,0)+\varepsilon>0.
\]
Then $\widetilde D(x,x')=\widetilde b+\sum_i(w_i x_i+w_i'x_i')$
is positive on the entire physical cube. Since $w_i-w_i'=a_i$,
substituting $x=z,x'=1-z$ gives $D(z)$ exactly. Lift the numerator by
placing its logical coefficients on the first rail and zero on the second
rail. The positive-affine-ratio realization theorem then gives one actual
atom. This proof also covers $r=0$.
\end{proof}

The lift does not provide a representation of arbitrary labels on the
physical cube. For example, a natural mask of the complementary-pair
subset is $\prod_i(x_i-x_i')$, but multiplying an arbitrary rational
score by this mask is an additional operation not supplied by the model.
No efficient global mask or covering compiler is asserted here.

\paragraph{Formal status.}
The combinatorial extraction proof and local decoder theorem above are
written proofs only at this checkpoint. The analytic prefix and suffix bounds
and the exact nonnegative cancellation formula, including both vanishing
identities, are Lean-checked in
\texttt{Results/TorusPrismDecoder.lean}. The public declarations are
\texttt{rationalPrefixKernel\_error},
\texttt{rationalSuffixKernel\_error},
\texttt{gatedRationalInterval\_eq},
\texttt{gatedRationalInterval\_nonneg},
\texttt{gatedRationalInterval\_zero\_row}, and
\texttt{gatedRationalInterval\_equal\_endpoints}, all in the
\texttt{HeadComplexity} namespace. They use only the standard logical axioms
\texttt{propext}, \texttt{Classical.choice}, and \texttt{Quot.sound}.
The dual-rail proposition, including an actual atom witness, is Lean-checked
in \texttt{Results/DualRailLift.lean}, most directly as
\texttt{HeadComplexity.exists\_fracAtom\_dualRail\_affineRatio}. Its axioms
are the same three standard logical axioms.
The exact atom realization used in the proof is the existing Lean theorem
\texttt{HeadComplexity.FracAtom.ofPositiveAffineRatio\_eval} in
\texttt{Atoms/PositiveAffineRatio.lean}. No full-cube universal head upper
bound follows from the statement proved here.
\end{document}
